\section{Introduction}
\label{sec:intro}

Patent drawings encode design information through contours, dashed boundaries,
sectional views, and annotations. IMPACT~\cite{shomee2024impact} makes design
patent figures and associated captions available for multimodal analysis, while
PatentLMM~\cite{shukla2025patentlmm} studies descriptions of technical figures.
Generating 3D assets from such drawings could support interactive inspection of
design records. The task is nevertheless ambiguous: visible lines do not uniquely
determine depth, hidden surfaces, or whether adjacent figures depict one object.

Progress in generative modeling offers a starting point. Denoising diffusion
models~\cite{ho2020ddpm,song2020ddim} and latent diffusion~\cite{rombach2022ldm}
provide powerful image priors; diffusion transformers~\cite{peebles2023dit} and
flow-based approaches~\cite{lipman2023flow,liu2023rectified} offer complementary
ways to model complex visual distributions. In 3D, DreamFusion and
Magic3D~\cite{poole2022dreamfusion,lin2023magic3d} transfer image priors through
rendering, whereas LRM, TripoSR, and InstantMesh~\cite{hong2023lrm,tochilkin2024triposr,xu2024instantmesh}
learn efficient reconstruction mappings. TRELLIS~\cite{xiang2025trellis} generates
structured 3D latents that support multiple output representations.

These capabilities do not establish robustness to sparse technical sketches.
A line drawing provides little of the shading and texture present in an ordinary
image, and background extraction may remove the very strokes that specify shape.
SketchDream~\cite{liu2024sketchdream} explicitly addresses sketch and text
conditioning, while Know3D~\cite{chen2026know3d} uses vision-language knowledge to
control unseen regions. Their conditioning mechanisms motivate treating visual
structure and semantic plausibility as separate aspects of patent reconstruction.
Our approach investigates an input-level intervention on a frozen image-conditioned
generator rather than adding a language-conditioned generation mechanism.

We implement a sketch adaptation stage that preserves detected linework and
constructs a soft distance-to-ink representation. A Gaussian-smoothed version
provides an additional controllable scale of spatial support. Four matched
configurations separate ordinary TRELLIS preprocessing, bypassed background
removal, distance-map conditioning, and distance-map conditioning with smoothing.
All configurations receive the same reviewed crop and seed. The distance map is
a deterministic proximity field, not a learned boundary probability or a filled
object silhouette.

Evaluation is equally consequential. CLIP~\cite{radford2021clip} provides semantic
features, SSIM~\cite{wang2004ssim} measures structural image similarity, and
LPIPS~\cite{zhang2018lpips} supplies a learned perceptual distance. None alone
establishes recovery of the true 3D design. Background agreement, pose mismatch,
and thresholded linework can affect drawing-to-render comparisons. We therefore
combine image metrics, explicit coverage reporting, mesh measurements, and
qualitative inspection. Aggregating first within drawings and then within patents
avoids giving disproportionate weight to patents containing many drawings.

The current evidence comprises an archive of 438 reconstructions generated from patent TIFFs from
50 patents, a verified subset of saved image calculations, and an implemented
adaptation experiment whose matched GPU results remain pending. We contribute
(1) a reproducible distance-based input adaptation with controlled preprocessing
ablations; (2) a paired evaluation protocol separating seed, drawing, and patent
aggregation; and (3) an archive analysis identifying practical limitations of
mask extraction and multi-figure inputs. Quantitative gains from adaptation and
comparisons with external models must be established by the pending experiments.

\section{Related Work}
\label{sec:related_work}

\subsection{Generative Priors and Conditioning}
DDPM~\cite{ho2020ddpm} formulates generation as learned denoising, and
DDIM~\cite{song2020ddim} provides an alternative sampling process with the same
training objective. Latent diffusion~\cite{rombach2022ldm} performs synthesis in
a compressed representation; DiT~\cite{peebles2023dit} studies transformer-based
diffusion architectures. Flow Matching~\cite{lipman2023flow} learns conditional
vector fields, while Rectified Flow~\cite{liu2023rectified} encourages straight
transport paths. These works explain the broader priors available to visual
synthesis, but are not additional components trained in our pipeline.

ControlNet~\cite{zhang2023controlnet} learns spatial conditioning branches for
pretrained image diffusion models. IP-Adapter~\cite{ye2023ipadapter} introduces
image-prompt conditioning through decoupled attention. Our intervention instead
changes the image representation supplied to an existing 3D generator without
learning an adapter. CLIP~\cite{radford2021clip} and DINOv2~\cite{oquab2023dinov2}
illustrate complementary language-supervised and self-supervised visual features;
semantic similarity in one feature space need not demonstrate geometric fidelity.

\subsection{Representations and 3D Generation}
NeRF~\cite{mildenhall2020nerf} represents scenes as continuous radiance fields,
whereas 3D Gaussian Splatting~\cite{kerbl2023gaussians} uses explicit Gaussian
primitives for efficient rendering. DreamFusion~\cite{poole2022dreamfusion}
optimizes a 3D representation using a pretrained diffusion prior, and
Magic3D~\cite{lin2023magic3d} adopts coarse-to-fine optimization culminating in a
textured mesh. ProlificDreamer~\cite{wang2023prolific} introduces variational score
distillation to address limitations of score-distillation-based generation.

Point-E~\cite{nichol2022pointe} generates point clouds, while
Shap-E~\cite{jun2023shape} generates parameters of implicit 3D functions.
LRM~\cite{hong2023lrm} predicts a radiance-field representation from an image;
TripoSR~\cite{tochilkin2024triposr} develops a fast feed-forward reconstruction
system. InstantMesh~\cite{xu2024instantmesh} combines multiview generation with
sparse-view reconstruction. TRELLIS~\cite{xiang2025trellis} uses structured
latents that can be decoded into meshes, Gaussians, or radiance fields; it is the
backbone used here. Differences in representation and export require common
rendering and accounting before comparing these systems quantitatively.

\subsection{Multiview and Sketch Control}
Zero-1-to-3~\cite{liu2023zero123} conditions novel-view synthesis on relative
camera changes. Zero123++~\cite{shi2023zero123pp} generates consistent multiview
images from one image, and SyncDreamer~\cite{liu2023syncdreamer} coordinates
multiview diffusion. Wonder3D~\cite{long2023wonder3d} jointly generates multiview
normal and color information. MVDream~\cite{shi2023mvdream} learns a multiview
prior from text and 3D data. These methods distinguish controlled view synthesis
from merely rendering a completed mesh at several evaluation angles.

SketchDream~\cite{liu2024sketchdream} combines sketches, text, depth guidance,
and multiview attention for generation and editing. Know3D~\cite{chen2026know3d}
uses latent hidden-state injection to transfer vision-language knowledge into
unseen-region generation. We discuss them as relevant comparison methods, but
have no independently measured benchmark results for either in the current
archive. Our implemented adaptation does not claim their conditioning mechanisms.

\subsection{Data and Reconstruction Assessment}
ShapeNet~\cite{chang2015shapenet}, Objaverse~\cite{deitke2023objaverse}, and
Objaverse-XL~\cite{deitke2023objaversexl} provide large collections of 3D assets.
In contrast, IMPACT~\cite{shomee2024impact} supplies patent figures and captions,
not the ground-truth meshes required for direct geometric reconstruction error.
PatentLMM~\cite{shukla2025patentlmm} further emphasizes technical-figure language
understanding, a related but distinct task.

SSIM~\cite{wang2004ssim}, LPIPS~\cite{zhang2018lpips}, and
CLIP~\cite{radford2021clip} measure different kinds of image agreement. Their
use with sparse drawings requires attention to framing, white backgrounds,
annotations, and pose. We retain these complementary measures and report mask
and mesh diagnostics separately, rather than interpreting semantic alignment as
a substitute for geometric ground truth.
